\section{讲座定位：多模态不是在 LLM 外面再挂一个视觉模块}

这节课的跨度很大：前半段回顾 BERT、GPT-3、ChatGPT 三个语言模型“时刻”，中段讨论 Transformer 架构、训练系统、alignment 与数据工程，后半段才进入 BLIP-2、LLaVA、CogVLM、CogAgent、图像扩散与视频生成。这样的安排不是离题，而是在回答一个贯穿全课的问题：\textbf{当模型能力从 text 扩展到 image、GUI 和 video 时，哪些来自 LLM 的规律仍然成立，哪些必须换表示、目标函数或系统路径？}

\subsection{来源修复与课堂边界}

旧稿把课堂日期写成 2026 年 4 月 4 日，并加入 monitoring dashboard、drift detection、automatic rollback、fairness governance 等大量课程中不存在的内容。本讲现以官方 CS25 V4 课程页、2024 年 5 月 9 日课堂、Stanford Online 官方录像 \texttt{cYfKQ6YG9Qo}、官方 \texttt{en-US} 人工字幕和 31 个录像恢复 teaching states 为课堂主干。官方页面和视频描述均未提供可验证的独立 slide deck，因此录像共享屏幕是视觉 source of truth。

\lecturefigure{slide-01-title.jpg}{官方课堂标题：From Large Language Models to Large Multi-modality Models}{Stanford Online 官方录像 00:02:31}

标题中的 ``Large Multi-modality Models'' 是 2024 年课堂用语。它不表示所有 modality 已被统一解决，而是强调模型研究已经从单一语言接口扩展到视觉理解、图像生成、GUI action 和视频。讲者所在团队的 CogVLM、CogAgent、CogView 与 CogVideo 是案例，不是整门课的唯一结论。

\begin{warningbox}{时间边界必须保留}
本讲展示的 benchmark、产品能力、下载量和“一两年内”的预测都属于 2024-05-09 的课堂快照。后来的模型版本、排行榜变化或产业结果不能倒灌进课堂事实；若用于拓展，只能显式标成 post-lecture evidence。
\end{warningbox}

\subsection{四段路线图与阅读问题}

课程先问 why、再问 how、接着问 what、最后问 where。读者不应把四部分看成并列综述：LLM 的 objective 和 scaling 解释为什么系统工程变得重要；系统与数据解释为什么 VLM 可以迅速复用语言底座；视觉表示的缺陷又解释为什么生成侧从 autoregression 转向 diffusion。

\lecturefigure{slide-02-outline.jpg}{课程路线图：历史、训练技术、视觉语言模型与研究方向}{Stanford Online 官方录像 00:03:29}

\begin{importantbox}{带着四个问题读完全课}
\begin{enumerate}
  \item \textbf{Objective：}不同 modality 应该共享 autoregressive objective，还是使用专门的 diffusion / contrastive objective？
  \item \textbf{Architecture：}投影层、Q-Former、vision experts 与 cross attention 分别解决什么接口问题？
  \item \textbf{Systems：}ZeRO（Zero Redundancy Optimizer，在 data-parallel ranks 间分片 optimizer/gradient/parameter states）、parallelism、long context 与高分辨率 token 如何改变可训练问题的边界？
  \item \textbf{Data：}当通用 web data 接近饱和，如何把 compute 转换成更高质量、可验证的数据？
\end{enumerate}
\end{importantbox}

\teachervoice{00:02:31--00:04:35，Ming Ding 提醒听众本讲会跨越多个并不熟悉的子领域。他的目标不是给每个模型做完整 survey，而是沿 Transformer 社区的时间线解释研究重心为什么移动。}

\subsection{证据账本：机制、结果、判断与预测}

同一张 slide 可能同时包含论文结构图、课堂判断和产品结果。为了避免“图上写了”被误读为普遍事实，本讲采用四层证据账本。

\begin{knowledgebox}{四层证据读法}
\begin{tabularx}{\textwidth}{L{0.18\textwidth} L{0.31\textwidth} X}
\toprule
层次 & 例子 & 正确读法 \\
\midrule
标准机制 & cross entropy、attention、DDPM、DPO & 给出公式、变量与假设 \\
具体系统 & ZeRO、Megatron、CogVLM、MM-DiT & 说明接口、资源路径与版本 \\
课堂结果 & benchmark table、web-agent trace、速度数字 & 只支持展示的设置和 protocol \\
讲者判断/预测 & loss 决定能力、video 将成热点 & 保留日期、语气与反例空间 \\
\bottomrule
\end{tabularx}
\end{knowledgebox}

\subsection{本章小结}

这是一堂“从 LLM 规律迁移到 multimodality”的课，不是 Cog 系列产品说明书。读者需要同时追踪 objective、architecture、systems 与 data，并始终区分标准机制、课堂实验、讲者判断和未来预测。

